% GENERATED FILE. Edit canonical lessons, then run npm run generate:books.
% canonical-source-hash: fnv1a64:72fd246c76c54b47
% canonical-lessons: TE-C28-madhyahnam-widened

\chapter{Afternoon}
\label{ch:te-afternoon}
\hlchaptermodality{28}

\begin{chapteropening}
\textbf{By the end of this chapter:} \emph{I can use Telugu's noon word to talk about the afternoon, and reach for the classical term when I need to be precise.}

Reuse \te{మధ్యాహ్నం} for the stretch between noon and evening, then contrast it with \te{అపరాహ్నం}, the more formal, more precise later-afternoon word.
\end{chapteropening}

\section[madhyāhnam]{\te{మధ్యాహ్నం} — the same \textquotedblleft{}noon\textquotedblright{} word, now also \textquotedblleft{}afternoon\textquotedblright{}}

\label{lesson:TE-C28-madhyahnam-widened}



\begin{quote}
You already know \te{మధ్యాహ్నం} means \textquotedblleft{}noon\textquotedblright{} — precisely, \textquotedblleft{}middle of the day.\textquotedblright{} This lesson's honest twist: the word doesn't stay that precise in modern use.
\end{quote}

\begin{culture}
\textbf{\te{మధ్యాహ్నం}} (\textbf{madhyāhnam}) is the exact word from Chapter 17, still literally \textquotedblleft{}middle\textquotedblright{} + \textquotedblleft{}day\textquotedblright{} — but a source fetched directly (not just a search summary) confirms it now covers \textquotedblleft{}\textbf{the time between noon and evening}\textquotedblright{} in modern usage — a genuine widened sense, not just the precise midday instant its etymology points to. This mirrors Hindi's \textbf{\dv{दोपहर}} (clearly, well- confirmed widened from \textquotedblleft{}noon\textquotedblright{} to \textquotedblleft{}afternoon\textquotedblright{}) and, more thinly Kannada's own \textbf{\te{మధ్యాహ్న}} (a hedged, single-source impression of the same kind of stretch).
\end{culture}

\begin{cousinweb}
Telugu also has \textbf{\te{అపరాహ్నం}} (\textbf{aparāhnam}), from the same classical \textbf{five-part Sanskrit day division} already found behind Kannada's \textbf{\te{అపరాహ్న}} — \emph{prāta/udaya} (sunrise), \emph{saṅgava}, \emph{madhyāhna} (noon), \emph{aparāhna} (\textquotedblleft{}later afternoon\textquotedblright{} specifically), \emph{sāyāhna} (evening). Be honest about the sourcing here, echoing the same hedge already applied to Kannada's \emph{aparāhna}: thin, informal evidence — a native-speaker forum discussion, general dictionary usage notes — points toward \textbf{\te{మధ్యాహ్నం} winning colloquially}, with \textbf{\te{అపరాహ్నం}} staying the more \textbf{formally precise} term. Real, but not corpus-linguistics-grade sourcing — treat it the same cautious way Kannada's own equivalent finding was treated, not as fully settled.
\end{cousinweb}

\subsection*{Your turn}
Say these aloud:

\begin{itemize}
\raggedright
  \item \textquotedblleft{}madhyāhnam\textquotedblright{} — the same word from Chapter 17, \textquotedblleft{}middle of the day\textquotedblright{}
  \item the confirmed widening — now also covers \textquotedblleft{}the time between noon and evening\textquotedblright{}
  \item \textquotedblleft{}aparāhnam\textquotedblright{} — the more precise classical word for later afternoon, staying the more formal term per thin but real evidence
\end{itemize}

\begin{tcolorbox}[breakable,colback=teal!4,colframe=teal!35!black,title={Before you move on}]
Is \te{మధ్యాహ్నం} a new word, or the same one from Chapter 17? (\textbf{The same word} — \textquotedblleft{}middle\textquotedblright{} + \textquotedblleft{}day.\textquotedblright{}) Does modern usage confirm \te{మధ్యాహ్నం} covers only the exact instant of noon? (\textbf{No} — a directly- fetched source confirms it now covers \textquotedblleft{}the time between noon and evening\textquotedblright{} too.) Is it confirmed whether \te{అపరాహ్నం} or widened \te{మధ్యాహ్నం} is more common today? (\textbf{Thin but real evidence} points to \te{మధ్యాహ్నం} winning colloquially, \te{అపరాహ్నం} staying formally precise — not fully settled, but not silent either.)
\end{tcolorbox}
